% Generated draft: requires mathematical and source review.
\documentclass[11pt]{article}
\usepackage[margin=25mm]{geometry}
\usepackage{fontspec}
\setmainfont{Latin Modern Roman}
\usepackage{amsmath,amssymb,mathtools,unicode-math}
\setmathfont{Latin Modern Math}
\usepackage{enumitem,needspace,xurl,hyperref}
\setlength{\parindent}{0pt}
\setlength{\parskip}{5pt}
\setlength{\emergencystretch}{3em}
\newcommand{\N}{\mathbb{N}}
\newcommand{\Z}{\mathbb{Z}}
\newcommand{\Q}{\mathbb{Q}}
\newcommand{\R}{\mathbb{R}}
\newcommand{\C}{\mathbb{C}}
\newcommand{\F}{\mathbb{F}}
\newcommand{\A}{\mathbb{A}}
\newcommand{\PP}{\mathbb{P}}
\newcommand{\E}{\mathbb{E}}
\newcommand{\eps}{\varepsilon}
\newcommand{\dd}{\,\mathrm{d}}
\begin{document}
\section*{UnsolvedMath ID: 30001739}
\textbf{Source problem code:} catalogue entry
\subsection{Definitions and mathematical statement}
Let \(E/F\) be a quadratic extension of local fields with nonidentity Galois automorphism \(\tau\), and let \(G=\operatorname{GL}_n(E)\). Apply \(\tau\) entrywise to matrices and put
\[
X=\{h\in M_n(E):{}^t\tau(h)=h,\ \det(h)\ne0\}.
\]
The natural right action is
\[
h\cdot g={}^t\tau(g)hg\qquad(h\in X,\ g\in G).
\]
For a representation \(\rho\) of \(\operatorname{GL}_m(E)\), Galois invariance means \(\rho^\tau\simeq\rho\), where \(\rho^\tau(g)=\rho(\tau(g))\).

The supplied record attributes the following assertion to Conjecture 4 of the cited report. Given a finite list \(\pi_1,\ldots,\pi_k\) of \(\tau\)-invariant imprimitive representations, if their product
\[
\pi=\pi_1\times\cdots\times\pi_k
\]
is an irreducible representation of \(G\), then
\[
\dim\operatorname{Hom}_G(X,\pi^*)=2^k.
\]
If the factors are representations of general linear groups of ranks \(n_1,\ldots,n_k\), the intended product must have total rank \(n_1+\cdots+n_k=n\).

\textit{Unresolved qualifications.} The supplied record does not define ``imprimitive,'' the representation category, the product \(\times\), or the dual \(\pi^*\). Moreover, \(X\) as defined above is a variety with a group action, not a specified linear representation. Thus \(\operatorname{Hom}_G(X,\pi^*)\) retains the source notation: a function space, distribution space, or other linear interpretation of \(X\), and the meaning of equivariant morphism, must be supplied before this is a fully specified multiplicity. The record also does not state whether the local fields must be nonarchimedean or satisfy characteristic restrictions. The displayed assertion is consequently a statement-only transcription with unresolved conventions, not a completed precise formulation.
\subsection{Short English statement}
The supplied conjecture predicts multiplicity \(2^k\) for an irreducible product of \(k\) Galois-invariant imprimitive representations in the representation associated with nonsingular Hermitian matrices. The meaning of imprimitive representation and the representation-theoretic conventions needed to define this multiplicity remain unspecified in the supplied record.
\subsection{Sources}
\begin{itemize}
\item Supplied local UnsolvedMath registry record, permanent ID \(30001739\), source problem code \texttt{catalogue entry}, title \emph{Multiplicity of Galois-Invariant Imprimitive Products}; collection: Oberwolfach Reports Open Problems. No individual proposer is supplied; the proposed year is 2011.
\item \emph{Automorphic Forms: New Directions} (2011), Oberwolfach Reports, pp.~719--779, DOI \texttt{10.4171/owr/2011/14}. Supplied source and literature reference, labeled ``original OWR report'': \url{https://doi.org/10.4171/owr/2011/14}. The supplied assessment identifies the assertion as Conjecture 4; no conjecture-specific page or individual author is supplied.
\item Literature assessment dated exactly \texttt{2026-08-21}: the supplied curation labels the problem open and states that the results cited in the report do not settle the general case. This is supplied dated curation, not independently verified here; its time-sensitive status and cited claims require independent verification.
\item The registry labels the statement \texttt{corrected\_verified} and reports that its mathematical content was checked against the official report, with notation, context, and question form normalized without changing the source-level claim. This is supplied verification metadata, not an independent verification in this draft.
\end{itemize}
\end{document}
